In this study, building off the training sample augmentation technique originally introduced in \citep{2024ApJ...967L...6M}, we have developed, implemented and validated a SOM-based data-augmentation framework to improve the accuracy and robustness of photometric-redshift estimation. Starting from two distinct mock catalogues with different galaxy SED prescriptions -- \texttt{OpenUniverse2024} as a proxy for realistic LSST observations and \texttt{CosmoDC2} for targeted augmentation -- we conducted \(501\) experiments to generate photometric and spectroscopic datasets under both Y1 and Y10 survey depths.

We first constructed the \texttt{Application} datasets by applying LSST-like photometric errors, shear-noise and the DESC \texttt{Gold} selection to \texttt{OpenUniverse2024}, then trained a SOM on these samples to compress the high-dimensional colour–magnitude space into a two-dimensional grid of cells. Realistic spectroscopic \texttt{Degradation} datasets were produced by sampling this photometric population with magnitude, redshift and colour cuts, plus a logistic function of spectroscopic success rates, and by merging multiple independent noise and sampling realisations to emulate heterogeneous survey compilation.

To compensate for under-sampling in the degraded spectroscopic datasets, we used the \texttt{CosmoDC2} catalogues with the same photometric pipeline to map each galaxy onto the pre-trained SOM. We applied cell-specific weights to harmonise their overall distributions in SOM space with the photometric \texttt{Application} datasets. We then employed adaptive augmentation by targeting galaxies in SOM cells unoccupied by the \texttt{Degradation} datasets, or beyond their spectroscopic magnitude or redshift limits, and set the target augmentation-to-degradation ratio to \(f_\ast(1+f_\ast)\), where \(f_\ast\) is the fraction of \texttt{Application} galaxies in SOM cells unoccupied by the \texttt{Degradation} dataset.

Using \texttt{FlexZBoost}, a conditional‐density machine‐learning estimator, we trained photometric‐redshift models on the \texttt{Combination} datasets, formed by merging the \texttt{Degradation} and \texttt{Augmentation} catalogues. We quantified point‐estimate accuracy via the median scaled bias, NMAD scatter, outlier fraction and catastrophic‐failure rate for both lens and source samples, and used the divergence loss of the normalised PIT histograms as an empirical indicator of conditional-PDF performance. The comparison of models trained with and without augmentation, shown alongside the in-sample \texttt{Combination} reference, demonstrates the robustness of our approach and the clear efficacy of SOM‐based data augmentation.

Across both Y1 and Y10 scenarios, the lens sample typically achieves mean absolute scaled biases \(\overline{\left|\delta_z\right|} \sim 0.005\), below the SRD-motivated comparison benchmarks even without augmentation—demonstrating the precision attainable for bright galaxies. For source galaxies at \(z_{\mathrm{true}}\lesssim1.5\), the point-estimate diagnostics likewise lie near or below the comparison benchmarks. Whilst performance declines at \(z_{\mathrm{true}}\gtrsim1.5\), training on the \texttt{Combination} datasets with SOM-based augmentation reduces the average bias and catastrophic failure rate by up to a factor of two, with markedly lower divergence losses providing complementary empirical support for improved conditional-PDF modelling. Importantly, the reduced scatter across independent realisations shows that augmentation renders photometric-redshift performance less sensitive to variations in the spectroscopic selection function.

These results demonstrate that SOM‐based data augmentation robustly identifies the SED space without representative spectroscopic samples and compensates for spectroscopic incompleteness and mitigates SED mismatches, yielding significantly more accurate and reliable photometric-redshift estimates. Such improvements are essential for LSST cosmological and galaxy-evolution analyses, enabling precision measurements of weak gravitational lensing, large-scale structure and the formation history of billions of galaxies across cosmic time.
